\chapter{C14.3b: the complete gamma--circuit join and its negative soft source pivot}
\label{component:pd6aedcfef8b0}
\noindent\textbf{Component }\path{pd6aedcfef8b0}\par
{\small\noindent Source: \path{R17/proofs/c1_small_eta/coupled_gamma_soft_obstruction.tex}\par}
{\small\noindent SHA256: \path{d6aedcfef8b038e63739f4fe9c07831f1f1b8e419259cf37a2ab4410611f73cc}\par}
{\small\noindent Inventory use: GS1 exact polarization only | C14\_\allowbreak{}ENTIRE\_\allowbreak{}ORIGINAL\_\allowbreak{}PARENT\_\allowbreak{}GRAPH; precise subordinate uses retained below\par}
\medskip
\section{C14.3b: the complete gamma--circuit join and its negative soft source pivot}

\paragraph{Status and unchanged domain.}
This is a diagnostic of the fully conducted frame. The damped DG/DS
certificates elsewhere in this revision bypass it on their admitted
domains; it creates no additional obligation beyond the current frontier.
This note does not admit C14.3b. It gives an exact simultaneous join of
its remaining full normal with both coupled-circuit rows, and identifies
an additional obstruction after the former order-one angular obstruction
has been removed. Retain $x,\eta>0$ infinitesimal, $a=1/\zeta$ a positive
bounded unit, both original gap openings, $\beta\ge0$, independent
infinitesimal testing angles $u,v$, and the normalized $W$-unit state
chart. No source/testing rate comparison or mixed inverse is used.
The completed R16 bounded-$W/Z,X/Z$ chart and literal $Z=0$ endpoint
are unchanged.

Use the homogeneous JF1 notation, and abbreviate $Q=Q_t$, $T=T_D$,
$G=\gamma h$, $x=k^2$, $a=\eta c^2$. Put
\[
 H=avW-uZ,\quad \tau=\eta cW-Z,\quad D_0=u-cv,
 \quad J=(uv-x/c)Z,
\]
\[
 E=\eta uW-vZ=\eta(P+Q),\quad
 R_c=kv\tau,\quad R=-\sqrt\lambda\,\eta kD_0W,
 \quad \mathcal T^2=(2\lambda x/c)uv\tau^2.
\]
The inherited gauges allow $u,v\ge0$; the identities below themselves
are polynomial and do not require a testing ratio inverse. In
particular no inverse of $Z$ is used.

\paragraph{Join the entire remaining normal before extracting the circuit.}
The exact conductor and polarization are
\[
 uv\xi=h\{uJ+kR_c-(x/c)H-T\},\qquad
 EH=uv\tau^2-\eta D_0^2WZ.                       \tag{GS1}
\]
Define the source-only shift
\[
 d=\frac{G}{2\lambda},\quad
 \widehat R=R+\frac{Gk}{2\sqrt\lambda}Q,
 \quad b_\gamma=b'-G/2,
\]
and the complete modified source packet
\[
 \begin{split}
 S_\gamma={}&CP^2+AQ^2+D_p(P-hT)^2+2b_\gamma QT\\
             &-\lambda x(E-dQ)^2.
 \end{split}                                                   \tag{GS2}
\]
Then the entire affected block has the exact identity
\[
\begin{split}
 &CP^2+AQ^2+D_p(P-hT)^2+2b'QT
       +R^2+\mathcal T^2-\lambda xE^2+\gamma uvQ\xi\\
 ={}&S_\gamma+\widehat R^2+GuQJ+(Gx/c)PH\\
 &+\frac{x}{c}(2\lambda-G/\eta)uv\tau^2
                    +(Gx/c)D_0^2WZ.                 
\end{split}\tag{GS3}
\]
This joins the full $R_c$ coupling, its $R$ companion, the negative
$E$ correction and the extra circuit row at once. The completion
\[
 R^2+(Gk/\sqrt\lambda)QR+GxQE-\lambda xE^2
           =\widehat R^2-\lambda x(E-dQ)^2
\]
is exact. The whole angular $EH$ term has canceled in GS3. Thus the
order-one negative angular residue JR7 is not being reused as the
obstruction in this note. All the remaining JF1 terms are left
unchanged, including $\lambda J^2$, the complete radial rows, the
angular block, the full $D_p$ shift and every higher-angle square.

\paragraph{The newly exposed source block is strictly indefinite.}
First omit just the negative last square in GS2, and call the result
$S_0$. Its pure-source $P$ completion is
\[
 S_0=(C+D_p)\left(P-\frac{hD_p}{C+D_p}T\right)^2
       +AQ^2+h^2K T^2+2b_\gamma QT,
 \qquad K=\frac{CD_p}{C+D_p}.                       \tag{GS4}
\]
On $(\sqrt\eta Q,T/\sqrt\eta)$ the remaining determinant is
\[
                 \Delta=h^2AK-b_\gamma^2.           \tag{GS5}
\]
There is a particularly simple exact native-source identity:
\[
 b_\gamma=-\frac{hC_n}{r},\qquad
 C_n=(1-r)C=\frac{(1+x)(r+\eta)E_s}{2P_d},
 \quad E_s=1-x\eta h.                              \tag{GS6}
\]
Indeed $b_\gamma=(h/r)(\eta B_x(1-h)-D_pj)$, and the original
formulas give $\eta B_x(1-h)-D_pj=-C_n$. This identity holds at
literal $\beta=0$ as well.

For a gap-free expression put
\[
 T_s=1+xh,\quad j=\eta T_s+r-1,\quad
 L=T_s-xhj+x(1-r)/a,\quad
 P_4=T_sE_s+x(1-r)(h+1/a).
\]
The unchanged source has $A=\eta P_4/(2rP_d)$ and
$K=C_n(1+y)/(\eta L)$, hence
\[
 \boxed{\quad
 \Delta=\frac{h^2C_n}{2r^2P_dL}
       \{r(1+y)P_4-(1+x)(r+\eta)E_sL\}.
 \quad}                                                     \tag{GS7}
\]
Here $r,h,L,P_d$ are positive source units. Neither original gap
appears in a denominator.

Write $h=1-x\pi$, $r=1-\eta\sigma$ with the inherited bounded
nonnegative $\pi,\sigma$. Exact finite source division, not a
comparison of $x$ and $\eta$, gives
\[
 \boxed{\quad \Delta=-x d_x-\eta d_\eta,\qquad
 d_x\sim\tfrac14,\quad
 d_\eta=\frac{r+\eta}{4r^2}\sim\tfrac14.
 \quad}                                                     \tag{GS8}
\]
For an explicit definition, set
$d_x=-\{\Delta-\Delta|_{x=0}\}/x$; the numerator vanishes identically
at $x=0$. Formula GS7 has a unit denominator there, so this is a
bounded source rational function. Also
$\Delta|_{x=0}=-\eta(r+\eta)/(4r^2)$ and
$\partial_x\Delta|_{x=\eta=0}=-1/4$, proving all of GS8 uniformly
through the gap openings. In particular $\Delta<0$ on the entire
unchanged source chart, regardless of $x/\eta$.

This is an actual independent-root negative direction. Choose
$Q=1$, $T=-b_\gamma/(h^2K)$ and
$P=hD_pT/(C+D_p)$. Then
\[
                      S_0=\Delta/(h^2K)<0.
\]
Subtracting the last nonnegative square in GS2 cannot repair it;
$S_\gamma<0$ at this same independent-root direction. Therefore
one cannot extract the two new circuit terms in GS3, complete
$\widehat R$ independently, and claim that the remaining source
block has a nonnegative, let alone strict, bounded Gram. This
obstruction survives the simultaneous $R$ completion, which was
not performed in the earlier failed angular-residue calculation.

\paragraph{The signed circuit remainder also persists.}
The other coefficient in GS3 has the exact finite source division
\[
 2\eta\lambda-G=-x d_c-\eta d_{\eta c},\qquad
 d_c\sim\pi+1+\frac1{2a}>0,\quad
 d_{\eta c}=\sigma(r+\eta)/r\ge0.          \tag{GS9}
\]
The quotient has only source-unit denominators. The numerator used
for the $x$ division vanishes at $x=0$ after subtracting the exact
value $-\eta\sigma(r+\eta)/r$. Consequently the last line of GS3 is
\[
 -\frac{x(xd_c+\eta d_{\eta c})}{\eta c}uv\tau^2
                   +\frac{Gx}{c}D_0^2WZ.           \tag{GS10}
\]
On the pure-testing sign chart $WZ\ge0$ the second term is a literal
positive square, but the first is negative. On $WZ<0$ that second
term changes sign too. Thus no sign of the state or source is being
discarded in the proposed simultaneous join.

\paragraph{Two exact unit-$W$ dependencies that a further join must preserve.}
Write $f_b=F_b/\eta=uvW-(x/a)Z$. Then
\[
 Zf_b-WJ=\frac{x}{a}Z\tau,
 \qquad
 \sqrt{x}\frac E{\sqrt\eta}
       =-\frac{R}{\sqrt{\lambda\eta}}
                  +\sqrt{\frac{x}{\eta}}v\tau.      \tag{GS11}
\]
The first $J$ in GS11 is the homogeneous root $(uv-x/c)Z$.
These identities use $W$, $Z$ multiplicatively and remain valid at
$Z=0$. They are relevant to the soft scale exposed in GS8, but do
not by themselves provide a strict independent $\sqrt{x\eta}Q$
pivot or authorize another bare radial square. In particular the
positive roots and state coordinates are dependent; the negative
independent-root direction above is not a counterexample to the
original parent.

At the original contact all roots in GS1--GS3 vanish, including
$P,Q,T,J,R,\tau$; equality and every previously retained normal are
preserved. The exact full-parent identity is obtained by adding
unchanged JF1 terms to GS3, so there is no dropped normal or packet.
The remaining obligation is now concrete: compensate both the
negative source scale $x+\eta$ in GS8 and the signed two-row
remainder GS10 using the actual circuit/radial/state dependencies.
No such compensation, and no all-tests admission, is asserted here.

The focused checker \path{check_coupled_gamma_soft.py} verifies
GS1--GS3, the source completion, the gap-free native-source reduction,
the two finite divisions and residues, and GS11. It does not rerun
inherited suites or compare source/testing rates numerically.
